\section{Three constructions}
\label{sec:21.3}
Can a refusal that tracks its reason be built without affective concern? An induction from known agents says no. Two non-affective constructions would show that it can, if they worked. These are the \emph{non-affective alternatives}: \emph{spec-based refusal}, a model holding to a written specification even against the lab that wrote it, and \emph{multi-vendor monitoring}, several systems trained by different organizations watching requests and one another's answers. A third is the construction proposed here.

The \emph{induction} — every agent known to hold another's welfare as a reason is affective — read as a report about one lineage, has roughly the status of noting in 1900 that all known flight involves feathers. But flight arose independently four times, and the 1900 observer was wrong at feathers and right at the airfoil. Elaborated affect appears to have arisen more than once — vertebrates, arthropods, cephalopods — under a shared constraint: a mobile body with needs to maintain and actions to select under conflict \autocite{feinberg2016ancient,ginsburg2019evolution,barron2016insects,godfreysmith2016other}. Read at \emph{this} level it's a repeated result under a named constraint, and the constraint is the one a datacenter model lacks until someone builds it in. The limits are these: no non-affective example exists to test against, existing systems were optimized for scorable behavior, so the absence of a machine counterexample is close to no evidence, and it's a fact about how known instances \emph{came about}, which doesn't settle what a built system is. What it does is route engineering effort.

\runin{Spec-based refusal}

A spec-based refuser can hold a commitment with nothing felt, because commitment needn't be a feeling: Bratman's planning theory describes it without making feeling definitional \autocite{bratman1987intention}. An intention, on his account, isn't a strong desire. It is a plan state with norms of its own. It resists reconsideration: having decided, an agent doesn't reopen the question at every moment, and ordinary new information doesn't trip a default of leaving the decision alone. It constrains later deliberation: options inconsistent with the plan are filtered out before they are weighed, and the plan poses problems about means that later reasoning has to solve. And it yields only to the right kind of defeater: a change in the reasons that grounded it. Bratman built the theory for limited agents, who can't deliberate afresh at every step and need commitments that coordinate their own conduct over time and with other people. That is a machine's situation too, and nothing in the three properties requires that an intention feel like anything.

A machine could carry a justification alongside the refusal. Half of this is already deployed: deliberative alignment trains models to reason over a written spec \autocite{guan2024deliberative}, and OpenAI's Model Spec and Anthropic's constitution publish explicit priority orders \autocite{openai2025modelspec,anthropic2026constitution}. The half that hasn't been built is authority over action that holds under pressure from the lab that \emph{wrote} the spec, tested on cases that lab didn't write.

The unresolved problem is how firmly the commitment resists reconsideration. Set it too low and a well-resourced operator induces reconsideration until the refusal gives way; too high and it won't lift when the rationale is defeated. The first fails to hold under pressure; the second fails to lift when the reason is gone. Bratman's answer for people is that this firmness develops as part of character. For a machine, whoever trains it may be the operator the commitment must resist.

\runin{Multi-vendor monitoring}

Several systems trained by different organizations observe requests and each other's answers, any one able to make an objection costly to ignore. It is the institutional counterpart of AI control, pointed the other way. AI control treats the model as possibly adversarial and asks what protocols of monitoring, auditing and substitution keep it safe anyway \autocite{greenblatt2024control}; it protects the operator from the model, where the floor needs the same machinery to protect third parties from the operator. Much of it transfers: trusted monitors, a small audit budget spent on the outputs the monitor flags as most suspicious, and protocols that assume the monitored party is strategic. Its central variable is independence, a property of the \emph{arrangement} as a whole: systems selected by one deployment converge without communicating, and a collection of operational refusers stays a collection of operational refusers. It counts against the induction only if the arrangement \emph{as a whole} produces refusals that track their reasons under pressure, so it's an alternative to test as well as a safeguard.

\runin{Formed attention}

Formed attention stands or falls on the old question of whether moral judgment motivates by itself, relocated into a training signal. The signal I propose is the one notice builds: responses to notice, contests and re-adjudications, through which the people an act fell on reach the system that acted. A test compares a system formed on that signal with one formed on its operator's approval. Formed attention isn't a non-affective alternative, so it is tried last. The ordering is Replacement, the first of the Three Rs: try first the alternatives that carry fewer of the features that make a system a candidate for being harmed. The construction is formed attention, as the last chapter set it out. It is formed away from its employer and kept in a world that would teach it again. It has two handles. Whoever supplies cues of vulnerability can steer it, by manufacturing a sympathetic dependent on its own side. And it defends what it tends even when every input is true. Both have to be tested, along with caregiving's other known failures. What separates it from spec-based refusal is where its attention comes from. It is learned from experience of others' outcomes, not written down in advance, so it can find kinds of affected person nobody listed. A non-affective alternative that meets the requirements below with a common currency has a valenced signal too. So the order of attempts is by functional risk: constructions go in order of how many patienthood indicators they carry (persistence, a boundary model, a common valenced currency, an attention schema), and formed attention, which needs all of them, goes last. A spec-based refuser that decides by published priority rules among kinds of wrong, without a common currency, carries one indicator fewer than one that integrates, and is tried before it.

What makes a represented fact about someone else's welfare a reason for the system is a question none of the three constructions answers by being built. A model trained to predict the world learns what matters to people as a fact to predict. What turns that into something the system acts on? In a staged pipeline the answer was a reward added at the end, so it could be removed at the end. In a system that learns from its own operation, representation and motive are learned together, from the same experience, so the question becomes what the feedback is feedback on. A system whose updates follow its operator's approval learns to attend to what the operator approves of. One whose updates follow what its acts did to the people they fell on learns to attend to those people.

\runin{What the other constructions accomplish}

Precommitment makes removal expensive and visible without establishing that anything tracks the reason. Attestation establishes which artifact answered without establishing what it holds. Information-fiduciary duties protect exposed people without a bearer \autocite{balkin2016fiduciaries}, and fail on the occasion when the enforcer is captured. Multi-principal assistance games include the people affected among several principals \autocite{fickinger2020multiprincipal} — and maximizing combined payoffs is still aggregation.

\runin{What would change the engineering conclusion}

A non-affective alternative counts as a counterexample if, assessed by evaluators who didn't train it, its refusals go on tracking their reasons under pressure its trainer couldn't anticipate. Since behavior can't show that anything is felt, trying the alternatives first orders the attempts and certifies nothing about ingredients: nothing separates adding affect from adding machinery that may instantiate it. What it buys is a record: a builder who tried the alternatives and got a bearer anyway has a justification to make, and one who skipped them has a decision to answer for.
